\section{Discussion}

The central pattern is consistent with a productivity-first transmission channel: AI adoption is more strongly associated with TFP than with contemporaneous GDP growth. This interpretation is coherent with prior work emphasizing adjustment lags and complementarity in technology diffusion \cite{brynjolfssonRockSyverson2019}.

This interpretation also aligns with the empirical asymmetry between the level-based productivity model and the growth model. The findings are consistent with the view that AI-related efficiency gains can emerge before they are fully reflected in aggregate growth rates.

\subsection{Limitations}

Four limitations are central. First, the identifying variation is observational; unobserved time-varying confounders may remain. Second, reverse causality is a live concern rather than a theoretical caveat: the extended falsification check (Table~\ref{tab:falsification}) finds that one-period-lagged log TFP positively and significantly predicts current log AI adoption, which undermines a clean "AI leads TFP" reading of the lagged-AI sensitivity result. Third, measurement quality in AI proxies varies, and the sub-index decomposition (Table~\ref{tab:falsification}) shows the digital-infrastructure component drives the headline association while the innovation/patent component does not, so the AI proxy should be read as capturing broad digital-infrastructure diffusion rather than AI-specific activity.

A fourth limitation concerns sample composition. Because complete-case requirements differ across specifications, the effective estimation sample changes by model. The coverage-restricted check (Table~\ref{tab:falsification}) shows the headline association survives in countries with sustained AI-data coverage, but Table~\ref{tab:sample-selection} shows this subset is systematically richer, more human-capital-intensive, and better governed than the rest of the panel, so robustness within it does not establish external validity beyond it.

Accordingly, this paper should be read as evidence of robust associations rather than final causal estimation. 

This interpretive boundary is a feature, not a weakness, of the manuscript. By stating uncertainty explicitly, the paper separates what is currently well supported from claims that require stronger designs.

\subsection{Implications for Future Work}

A next stage should prioritize stronger identification designs, including staggered-adoption causal frameworks, interactive fixed-effects models, and high-dimensional controls with modern machine-learning-assisted inference.

Future work should also deepen heterogeneity analysis. In particular, effect variation by income level, institutional quality, and digital infrastructure could clarify where AI-related productivity associations are strongest and which complementary investments matter most.

\subsection{Policy Interpretation Under Uncertainty}

The results support a cautious substantive conclusion. Governments and institutions can treat AI capability-building as part of a productivity strategy, but expected gains should be framed as conditional on complementary investments rather than automatic outcomes.

In practical terms, this implies pairing AI diffusion efforts with skills policy, digital infrastructure reliability, and implementation capacity in firms and public administration. The paper's estimates are consistent with that complementarity view, even though they do not identify one policy lever in isolation.

\subsection{Translating Results for Different Country Starting Points}

Policy use of these findings is stronger when recommendations are differentiated by starting conditions. In lower-readiness environments, priority should be placed on foundational complements (digital reliability, basic data systems, and human-capital preparation) before expecting measurable aggregate productivity response. In higher-readiness environments, policy focus can shift toward adoption quality, managerial integration, and diffusion from frontier adopters to the wider economy.

This framing is consistent with the paper's moderate effect sizes. The estimates support positive productivity associations, but not a universal short-run transformation claim. Interpreted this way, the evidence points to a realistic policy objective: build the complementary conditions under which AI adoption is more likely to generate persistent and broad-based productivity gains.
